\apendice{Plan de Proyecto Software}

\section{Introducción}
El plan de proyecto software es el documento que define cómo se va a gestionar y ejecutar el desarrollo de un producto o sistema informático. Establece el alcance, los objetivos, el cronograma, los recursos, el presupuesto, los riesgos y los criterios de calidad, así como la organización del equipo y los mecanismos de seguimiento y control. En esencia, convierte una idea de software en un itinerario gestionable, con responsabilidades claras y procedimientos para controlar el progreso y los cambios.

Su importancia radica en que reduce la incertidumbre y aumenta la probabilidad de éxito académico y técnico. Permite justificar decisiones, anticipar problemas, y demostrar competencias en gestión de proyectos, ingeniería de requisitos, calidad y comunicación técnica.

En este documento se va a incluir la planificación temporal de las fases y tareas del proyecto. Para cada una se debe establecer una fecha de inicio y se debe estimar la fecha final. Para realizar esto hay que tener en cuenta el peso de cada tarea, los requisitos previos y las dependencias.

También, se va a valorar la viabilidad del proyecto. En concreto, la viabilidad económica y la legal. En la viabilidad económica se va a estimar los costes y los beneficios del proyecto. Mientras que en la viabilidad legal se van a valorar las leyes y licencias que afecten al desarrollo del proyecto.  

\section{Planificación temporal}
Se va a utilizar una metodología Scrum. En la cual se van a aplicar iteraciones (\textit{sprints}). Al finalizar cada \textit{sprint} se realiza una reunión para la revisión del incremento y la planificación del siguiente. En la planificación se deciden y valoran las tareas (\textit{issues}) a realizar en el siguiente \textit{sprint}. 

Estas tareas se han valorado usando los story points de ZenHub. 

\subsection{Sprint 0 (12/09/2025 - 19/09/2025)}


\subsection{Sprint 1 (19/09/2025 - 26/09/2025)}

\subsection{Sprint 2 (26/09/2025 - 03/10/2025)}

\section{Estudio de viabilidad}

\subsection{Viabilidad económica}

\subsection{Viabilidad legal}


